\documentclass[10pt,letterpaper,sans]{moderncv}
\moderncvstyle{banking}
\moderncvcolor{blue}

\usepackage[scale=0.84]{geometry}
\nopagenumbers{}
\setlength{\hintscolumnwidth}{2.2cm}

\name{Srikrishnaa}{Vadivel}
\title{Machine Learning Engineer | Scientific ML | Equivariant Models}
\email{srikrishnaacareers@gmail.com}
\phone[mobile]{516-853-5663}
\social[github]{krishnaa423}
\extrainfo{\href{https://leetcode.com/u/krishnaa42342}{leetcode.com/u/krishnaa42342} \textbar{} \href{https://hub.docker.com/u/krishnaa42342}{hub.docker.com/u/krishnaa42342}}

\begin{document}
\makecvtitle

\section{Profile}
\cvitem{}{PhD researcher applying machine learning to first-principles materials and molecular systems. Experience spans PyTorch, equivariant graph networks, diffusion models, distributed training concepts, HPC data pipelines, and physics-based loss design.}

\section{ML Experience}
\cventry{2021--present}{PhD Researcher}{Yale University}{}{}{\begin{itemize}%
\item Developed first-principles datasets and workflows for self-trapped exciton calculations on leadership-class HPC resources.
\item Prototyped EGNN-style force models with energy gradients $F_{ai}=-\partial E/\partial R_{ai}$ and molecular graph construction.
\item Built diffusion-model notes and code scaffolds covering DDPM training, conditional inference, transformers, and index-notation derivations.
\end{itemize}}
\cventry{2019--2021}{Software and Embedded Engineer}{Probe Technologies / Weatherford}{}{}{\begin{itemize}%
\item Built GPU-accelerated processing and visualization for large acoustic waveform data streams using CUDA and OpenGL.
\item Engineered production data tooling around ASP services and Microsoft database systems.
\end{itemize}}

\section{Selected Projects}
\cvitem{EGNN QM9 Force Model}{PyTorch message-passing model, force loss, synthetic graph fallback, and training-curve generation.}
\cvitem{Diffusion Notes}{DDPM equations with $\alpha_t$, $\beta_t$, conditional training, and inference in index notation.}
\cvitem{LiF Band Scaffold}{Reciprocal-space Hamiltonian assembly and eigenvalue plotting for physics-informed ML context.}

\section{Education}
\cventry{2021--present}{PhD, Electrical Engineering}{Yale University}{}{}{}
\cventry{2023}{MPhil, Electrical Engineering}{Yale University}{}{}{}
\cventry{2023}{MS, Electrical Engineering}{Yale University}{}{}{}
\cventry{2017--2018}{MEng, Electrical and Computer Engineering}{Cornell University}{}{}{}
\cventry{2013--2017}{BA, Electrical and Computer Engineering}{Cornell University}{}{}{}

\section{Skills}
\cvitem{ML}{PyTorch, scikit-learn, Hugging Face Transformers, Accelerate, datasets}
\cvitem{Scientific}{NumPy, SciPy, pandas, PETSc, SLEPc}
\cvitem{Systems}{Python, C++, Linux, Docker, Kubernetes, MPI, OpenMP}
\cvitem{Compute}{CUDA}

\end{document}
